\section{The finite approximation problem}\label{b:introduction}
A cylinder in $\R^3$ is a set $B+\R u$, where $u$ is a Euclidean unit
vector and $B\subset u^\perp$ is measurable with finite area. Its cost is
$|B|$; the cost of a finite family is the sum of these areas, counted with
multiplicity. A continuously varying family of segments suggests a way to
choose economical directions. Turning it into finitely many cylinders is
a separate problem: simply freezing a direction on each small label patch
can open gaps or create an area loss that does not vanish with the mesh.

The motivation comes from cylinder coverings of convex bodies. If
$K\subset\R^3$ is a compact convex set with nonempty interior, write
$A_{\min}(K)=\min_{|u|=1}|\pi_{u^\perp}K|$ for its smallest orthogonal
projection area. Bezdek records a question attributed to Bang asking
whether every finite cylinder cover costs at least $A_{\min}(K)/2$, together
with a two-cylinder equality example for the regular tetrahedron
\cite[Problem~3.1]{Bezdek2009}. Bezdek and Litvak prove the universal
lower bound $A_{\min}(K)/3$ in dimension three as a consequence of their
stronger directionwise normalized inequality
\cite[Theorem~3.1]{BL2009}. A useful way to search for inexpensive covers is
to perturb the lines in a simple configuration, compare the added label
area with the saving in perpendicular projection, and only then make the
family finite. This paper supplies that last step with its exact
hypotheses and area factor.

\subsection{Segments, tiles, and their Euclidean cost}
We use an axial coordinate $s\in\R$ and transverse coordinates
$p=(p_1,p_2)\in\R^2$. For a compact set $D\subset\R^2$, a continuous
field $V$ defined near $D$, and $L>0$, set
\begin{equation}\label{b:ruled-set}
 E(D,V,L)=\{(s,p+sV(p)):p\in D,\ |s|\le L\}.
\end{equation}
The label $p$ is the intercept at $s=0$ and $(1,V(p))$ is the
segment direction. A measurable tile $T\subset\R^2$ and a constant
velocity $g\in\R^2$ determine
\[
 \Cyl(T,g)=\{(s,\xi):s\in\R,\ \xi\in T+sg\}.
\]
At each height its section is a translate of $T$, rather than a deformed
copy. These are the cylinders used in our finite covers.

For a fixed $h>0$, allow the physical coordinate map
\[
 F_h(s,\xi)=(s,\xi_1,h\xi_2),\qquad
 J_h(g)=(1+g_1^2+h^2g_2^2)^{-1/2}.
\]
When $h=1$ the coordinates are orthonormal.
Throughout, a velocity's components in $J_h$ refer to these original
coordinates, even when the tile is rotated.

\begin{proposition}[Perpendicular base area]\label{b:projection}
Let $h>0$, let $T\subset\R^2$ be bounded and measurable, and let
$g\in\R^2$. The image $F_h(\Cyl(T,g))$ is a cylinder whose perpendicular
base has area $h|T|J_h(g)$. A compact nondegenerate parallelogram tile
has a compact nondegenerate parallelogram base. An open tile gives an
open base. The formula is unchanged after a Euclidean isometry.
\end{proposition}
\begin{proof}
The intercept tile $F_h(\{0\}\times T)$ has physical area $h|T|$ and
lies in the plane with unit normal $(1,0,0)$. Its cylinder direction is
$v=(1,g_1,hg_2)$. Orthogonal projection onto $v^\perp$ has constant area
Jacobian $|(1,0,0)\cdot v/|v||=J_h(g)$ on the intercept plane: this follows
by projecting an orthonormal pair and taking the area of its parallelogram.
A point and its projection differ by a multiple of $v$, so the projected
tile sweeps precisely the same cylinder. The projection is an invertible
linear map between the two planes because the first component of $v$ is
nonzero. It preserves compactness, openness, and nondegeneracy as asserted.
Isometries preserve perpendicular area and the cylinder property.
\end{proof}

Thus the prospective cost of the physical family $F_h(E(D,V,L))$ is
$h\int_D J_h(V(p))\,dp$. The following theorem realizes this upper
bound to arbitrary accuracy with square intercept tiles.

\begin{theorem}[Finite approximation by square tiles]\label{b:approximation}
Let $D\subset\R^2$ be compact, let $V$ be $C^1$ on an open
neighborhood of $D$, and let $L,h>0$. Suppose that
\begin{equation}\label{b:main-condition}
 \tr DV(p)=0,\qquad -L^{-2}<\det DV(p)\le0
 \quad(p\in D).
\end{equation}
For every $\varepsilon>0$, there are finitely many compact
nondegenerate square tiles $T_i\subset\R^2$, with possibly different
orientations, and velocities $g_i\in\R^2$ such that
\begin{equation}\label{b:main-cover}
 E(D,V,L)\subset\bigcup_i\Cyl(T_i,g_i),\qquad
 \sum_i |T_i|J_h(g_i)
 \le\int_D J_h(V(p))\,dp+\varepsilon.
\end{equation}
Under $F_h$, the perpendicular bases are compact nondegenerate
parallelograms and have total area $h\sum_i|T_i|J_h(g_i)$.
No condition on the area of $\partial D$ or smallness of $\|DV\|$
is required.
\end{theorem}

To obtain a prescribed error in physical base area, use that error
divided by $h$ as $\varepsilon$ in \eqref{b:main-cover}. The construction
includes compact label sets of area zero or with positive-area boundary.
The velocities $g_i$ are chosen from local affine approximations to $V$;
they need not equal the original field at the tile centers.

The eigenvalues in \eqref{b:main-condition} are
$\lambda_p,-\lambda_p$, where $\lambda_p=\sqrt{-\det DV(p)}$ and
$0\le L\lambda_p<1$. Thus opposite nonzero real eigenvalues are
allowed up to the strict length bound. When $\lambda_p=0$,
Cayley--Hamilton gives $(DV(p))^2=0$. In particular, a field with
square-zero differential everywhere on $D$ satisfies the theorem for
every fixed $L>0$, with no bound on its shear coefficient. The two types
may occur in the same field: for
\[
 V(p_1,p_2)=(p_2,p_1^3/3),\qquad
 DV(p)=\begin{pmatrix}0&1\\p_1^2&0\end{pmatrix},
\]
the theorem applies on any compact subset of $|p_1|<L^{-1}$.
The differential is square-zero on $p_1=0$ and hyperbolic elsewhere.

\subsection{Why fixed squares suffice}
For an affine field with derivative $A$, label differences at height
$s$ are transformed by $I+sA$. A trace-zero planar matrix has an
orthonormal basis in which its diagonal vanishes. In those coordinates,
write $A=\left(\begin{smallmatrix}0&b\\c&0\end{smallmatrix}\right)$.
The unit-grid centers then move to
\[
 (i+sbj,\,sci+j),\qquad(i,j)\in\Z^2.
\]
The squares themselves keep their sides and their area. They still
cover the entire plane whenever $0\le s^2bc<1$, even if an individual
coefficient $sb$ or $sc$ is large. A direct rounded fixed-point argument
proves this covering in Section~\ref{b:linear-coverings}. When one
coefficient is zero, the motion simply slides the rows or columns.

The bounded-parameter square covering is related to Stein's
notched-square construction in Kolountzakis's formulation
\cite[Theorem~7]{Kolountzakis1998}. The full positive-product range also
follows from the signed-function identity in the same paper
\cite[Section~2.3, equations~(17)--(20)]{Kolountzakis1998}.
We give an elementary proof, including the zero-product cases and all
square boundaries. Its integer iteration is a finite-chain instance of
Tarski's monotone fixed-point principle~\cite[Theorem~1]{Tarski1955}.
These comparisons concern the affine covering; the finite nonlinear
construction and its area estimate are proved here.

Section~\ref{b:finite-proof} passes to a $C^1$ field. On each coarse
square meeting $D$, we freeze an affine approximation and use a much
finer square grid in its orthonormal coordinates. The infinite affine
cover contains every actual nonlinear target. The inverse of $I+sA$
then shows that only centers in a thin collar of the coarse square are
needed, uniformly over the allowed heights. The fine-grid spacing and
the Taylor error both make that collar thinner than the coarse side
length by a factor tending to zero. Its extra tile area is therefore
negligible. Summing the sampled weights over the coarse squares gives
the integral over a shrinking outer neighborhood of $D$, which converges
to the integral over $D$ without a boundary-regularity assumption.

Section~\ref{b:product-section} develops a different construction with
open intercept cells bounded by hyperbolas. Their long arms absorb
expansion in one eigendirection and contraction in the other at small
relative area cost. This method packs the interior and requires a label
set with area-zero boundary; logarithmic dilation gives a finite-stopping
variant of the same construction. These alternatives are not needed for
Theorem~\ref{b:approximation}. Section~\ref{b:comparison} discusses
physical rescaling, the role of the spectral bound, and the use of a
finite approximation within a later parameter limit.
